\scanpage{63}
% manuscript line 1
% Exposed upper edge of the underlying sheet. The slip covers the lower strokes.
\[
\boxed{(a_1\wedge\cdots\wedge a_n,b_1\wedge\cdots\wedge b_n)=\det\big((a_i,b_j)\unsure{,}\big.} \tag*{(4)}
\]

% manuscript line 2
% Beginning of the overlying slip.
Soit $(e_i)_{i\in I}$ une base hilbertienne \inserted{tot. ordonnée} de \deleted{\unsure{$F$}} $E$, pour \deleted{toute} partie
finie \inserted{de $n$ éléments} \deleted{$J\subset I$} $J$ de $I$, posons

% manuscript line 4
\[
\boxed{e_J=e_{i_1}\wedge\cdots\wedge e_{i_n}}. \tag*{(5)}
\]

% manuscript line 5
\inserted{$J=(i_1,\ldots i_n)$, $i_1<i_2<\cdots<i_n$,} On voit aussitôt sur (4) que
$(e_J,e_{J'})=\delta_{J,J'}$ (symbole de Kronecker), donc les
$e_J$ forment un système orthonormal dans $\overset{n}{\bigwedge}\vphantom{\bigwedge}^{(2)}E$. Ce système
est \deleted{évidemment} même \uline{total}, \inserted{car l'espace $E_0$ engendré par les $e_i$ est dense dans $E$, et} \deleted{car l'espace \unsure{engendré}} l'espace vectoriel engendré
\deleted{de $E_0^n$ dans $\overset{n}{\bigwedge}E$ par l'application multilinéaire}
par les $e_J$ \inserted{dans $\overset{n}{\bigwedge}\vphantom{\bigwedge}^{(2)}E$} est identique à celui engendré par l'image
de $(E_0)^n$ par l'application multilinéaire naturelle de $E^n$
\deleted{dans} sur $\overset{n}{\bigwedge}E$, donc elle est dense dans $\overset{n}{\bigwedge}E$ puisque
$E_0$ est dense dans $E$, donc dense dans $\overset{n}{\bigwedge}\vphantom{\bigwedge}^{(2)}E$.

% manuscript line 14
% archive sheet number: 64
% Exposed lower portion of the supporting sheet.

% manuscript line 15
\unsure{normale} $(e_i)$, \deleted{telle que \unsure{l'on ait} \illegible{2}}

% manuscript line 16
\deleted{vérifie aussitôt que} \inserted{posons} $a_{\deleted{\unsure{\ldots}}\inserted{\textcircled{i}}}=\sum_j\deleted{\unsure{a_{ij}}}a_{ij}e_j$,
\deleted{\unsure{soit}} $A=(a_{ij})_{i,j\leq n}$, alors $(a_i,a_j)=\sum_k a_{ik}\overline{a_{jk}}$
\deleted{$\det((a_i,a_j))=\det$}
est l'élément d'indice $i,j$ de \inserted{la matrice} $AA^*$, donc
$\det((a_i,a_j))=\det AA^*=\det A\det A^*=\inserted{\det A\,\overline{\det A}=}|\det A|^2$, \deleted{d'où} il
reste à prouver \fbox{(6) $\lvert\det A\rvert\leq\|a_1\|\cdots\|a_n\|$} (inégalité
de Hadamard \unsure{!|?}), qui exprime que le volume du
parallélépipède \unsure{rectangle|construit} sur les $a_i$ est au plus
égal au produit des longueurs des arêtes).

% manuscript line 18a
% slanted
\margin{A\unsure{jouter} \unsure{préambule} \illegible{1} que $X\to x\wedge X$ de $\overset{n}{\bigwedge}\vphantom{\bigwedge}^{(2)}E$ dans $\overset{n+1}{\bigwedge}\vphantom{\bigwedge}^{(2)}E$ est de norme $\leq\|x\|$.}

% manuscript line 25
Pour ceci, \inserted{il suffit de remarquer qu'}on peut supposer les $(e_i)$ choisis de façon
que $A$ soit une matrice triangulaire, donc
$\det A=a_{11}\,a_{22}\cdots a_{nn}$, ce qui achève la démonstration.

% manuscript line 28
\uline{Corollaire 1} Soient $a_i,b_i$ $(i=1,\ldots,n)$ des éléments d'un espace
de Hilbert $E$. Alors on a $|\det((a_i,b_j))|\leq\prod_i\|a_i\|\,\|b_i\|$

% manuscript line 30
% archive sheet number: 63
